\chapter{Interface utilisateur}
\label{chap:interface}

\section{Design glassmorphique}

L'interface web adopte un style glassmorphique : fond sombre, éléments translucides avec \code{backdrop-filter: blur()}, typographie claire.

\section{Zone de saisie et résultats}

L'utilisateur saisit sa requête. La touche Entrée déclenche la recherche. Les résultats affichent : variables liées, score de pertinence, et preuves (bouton toggle).

\section{Panneau de diagnostic}

L'interface affiche : statut, warnings, joinStats (candidats testés, couples trouvés), pagination (pages et relations récupérées), AST.

\section{Gestion des erreurs}

Quand \code{statut: "Erreur"}, bannière rouge avec message. Aucun \code{undefined} affiché.
